% Generated draft: requires mathematical and source review.
\documentclass[11pt]{article}
\usepackage[margin=25mm]{geometry}
\usepackage{fontspec}
\setmainfont{Latin Modern Roman}
\usepackage{amsmath,amssymb,mathtools,unicode-math}
\setmathfont{Latin Modern Math}
\usepackage{enumitem,needspace,xurl,hyperref}
\setlength{\parindent}{0pt}
\setlength{\parskip}{5pt}
\setlength{\emergencystretch}{3em}
\newcommand{\N}{\mathbb{N}}
\newcommand{\Z}{\mathbb{Z}}
\newcommand{\Q}{\mathbb{Q}}
\newcommand{\R}{\mathbb{R}}
\newcommand{\C}{\mathbb{C}}
\newcommand{\F}{\mathbb{F}}
\newcommand{\A}{\mathbb{A}}
\newcommand{\PP}{\mathbb{P}}
\newcommand{\E}{\mathbb{E}}
\newcommand{\eps}{\varepsilon}
\newcommand{\dd}{\,\mathrm{d}}
\begin{document}
\section*{UnsolvedMath ID: 5300021}
\textbf{Source problem code:} catalogue entry
\subsection{Definitions and mathematical statement}
Let $\mathcal P$ be a finite collection of critically periodic quadratic polynomials $f(z)=z^2+c$. For a sequence of choices $f_i\in\mathcal P$, write $m_i$ for the exact period of its critical point $0$, so $f_i^{m_i}(0)=0$ and no smaller positive iterate returns to $0$. Retain the source's tuning operation and write
\[
g_n=f_1\vdash\cdots\vdash f_n,\qquad
n_k=\prod_{j=1}^k m_j.
\]
Here a superscript on $g_n$ denotes iteration, not pointwise exponentiation. For $k<n$ and the source's indices $i\leq n_k$, consider
\[
S_{n,k,i}=\left\{g_n^{\,n_k\ell+i}(0):\ell\in\Z,\ 0\leq\ell<m_{k+1}\right\}.
\]
The assertion is that these sets have geometry bounded uniformly as $n$, $k<n$, and $i\leq n_k$ vary, with the finite collection fixed. A bound may depend on that fixed collection, but should not deteriorate with tuning depth or the displayed indices. Whether it must also be uniform over all sequences drawn from the collection is not explicit.

The record does not define bounded geometry or its scale-invariant numerical criterion. Possible separation/diameter ratios, annular-modulus bounds, and other a priori geometry conditions cannot be identified without the source. It also does not define or delimit the tuning operation, give a lower bound for $i$, or specify whether $k=0$ is included. Nonnegative iterates require an appropriate nonnegative index range; the customary alternatives $1\leq i\leq n_k$ and $0\leq i\leq n_k$ are not silently equated. Degenerate critical periods and singleton or repeated-point sets require conventions for any geometry criterion. Thus the orbit sets and intended uniformity are specified, but the exact bounded-geometry theorem is missing a defining predicate.
\subsection{Short English statement}
Iterated tunings from a fixed finite list of critically periodic quadratics should produce certain small critical-orbit configurations with uniformly controlled geometry at every tuning depth. The source does not give the geometry criterion or complete index and tuning conventions.
\subsection{Sources}
\begin{itemize}
\item Local UnsolvedMath registry, permanent ID 5300021, exact source code catalogue entry, ``Uniform geometry in complex renormalization''; proposed by Mary Rees, 1992. Registry status open, classification OPEN-TRIAGE.
\item \emph{Stony Brook---Open Problems in Dynamical Systems}, 1992, Rees, \url{https://www.math.stonybrook.edu/open-problems-dynamical-systems}. Supplied extraction: source-TeX and pdf-text, accessed July 29, 2026. No PDF was accessed here.
\item The supplied research summary associates this uniform a priori geometry question with Rees's ``views of parameter space'' programme and with bounds of Lyubich and Kahn--Lyubich. No exact theorem or publication reference is supplied. It reports no verified general published proof of the stated claim; connections to a priori bounds are not treated as a resolution.
\item Provenance flags status NEEDS\_REVIEW: the source presents the item as open and source-side partial results were retained, but no comprehensive modern check was available. Public preprint access and provenance were verified; redistribution permission was not established and rights remain NEEDS\_REVIEW.
\end{itemize}
\end{document}
